\documentclass[12pt]{article}
\usepackage[spanish]{babel}
\usepackage[utf8]{inputenc}
\usepackage[T1]{fontenc}
\usepackage{amsmath, amssymb}
\usepackage{geometry}
\usepackage{enumitem}
\usepackage{needspace}
\usepackage{tikz}
\usetikzlibrary{babel}

\geometry{letterpaper, margin=2cm}
\setlength{\parindent}{0pt}
\setlength{\parskip}{6pt}
\hyphenpenalty=10000
\exhyphenpenalty=10000

\newcommand{\opcion}[2]{\item[#1)] $#2$}
\newcommand{\puntografico}{0.09}
\newcommand{\rotuloxeste}{\mbox{\textit{x} (este)}}
\newcommand{\rotuloynorte}{\mbox{\textit{y} (norte)}}

% Gráficas de circunferencia para opciones del ítem 4
% #1,#2 = centro; punto superior = (#1, #2+3)
\newcommand{\grafcirc}[2]{%
\pgfmathtruncatemacro{\topY}{#2+3}%
\begin{tikzpicture}[scale=0.42]
  \path[use as bounding box] (-7.5,-8) rectangle (10.5,9);
  \draw[<->, thick] (-6.5,0) -- (6.5,0);
  \draw[<->, thick] (0,-7.5) -- (0,7.5);
  \node[anchor=west, inner sep=1pt] at (6.55,0) {\scriptsize\rotuloxeste};
  \node[anchor=south, inner sep=1pt] at (0,7.85) {\scriptsize\rotuloynorte};
  % Punteadas: eje x-centro, eje y-centro, centro-punto superior, eje y-punto superior
  \draw[dashed, black!70] (#1,0) -- (#1,#2);
  \draw[dashed, black!70] (0,#2) -- (#1,#2);
  \draw[dashed, black!70] (#1,#2) -- (#1,\topY);
  \draw[dashed, black!70] (0,\topY) -- (#1,\topY);
  % Marcas y etiquetas (fill=white para visibilidad sobre punteadas)
  \draw (#1,-0.5) -- (#1,0.5);
  \node[fill=white, anchor=north, inner sep=1.5pt] at (#1,0) {\scriptsize $#1$};
  \draw (-0.5,#2) -- (0.5,#2);
  \node[fill=white, anchor=east, inner sep=1.5pt] at (0,#2) {\scriptsize $#2$};
  \draw (-0.5,\topY) -- (0.5,\topY);
  \node[fill=white, anchor=east, inner sep=1.5pt] at (0,\topY) {\scriptsize $\topY$};
  % Circunferencia y puntos
  \draw[thick] (#1,#2) circle (3);
  \fill (#1,#2) circle (0.18);
  \fill (#1,\topY) circle (0.18);
\end{tikzpicture}}

% Mini-gráficas para opciones gráficas de traslación
% #1,#2 = centro; #3 = radio
\newcommand{\grafcirctrasl}[3]{%
\pgfmathtruncatemacro{\rxcoordint}{#1+#3}%
\begin{tikzpicture}[scale=0.25]
  \path[use as bounding box] (-8,-15) rectangle (19.5,13.5);
  \draw[<->, thick] (-7,0) -- (14,0);
  \draw[<->, thick] (0,-14) -- (0,12);
  \node[anchor=west, inner sep=1pt] at (14.1,0) {\tiny\rotuloxeste};
  \node[anchor=south, inner sep=1pt] at (0,12.1) {\tiny\rotuloynorte};
  \draw[thick] (#1,#2) circle (#3);
  \draw[dashed] (#1,0) -- (#1,#2);
  \draw[dashed] (\rxcoordint,0) -- (\rxcoordint,#2);
  \draw[dashed] (0,#2) -- (\rxcoordint,#2);
  \fill (#1,#2) circle (0.42);
  \fill (\rxcoordint,#2) circle (0.42);
  \draw (#1,-0.45) -- (#1,0.45);
  \node[anchor=north, inner sep=1.5pt] at (#1,0) {\scriptsize $#1$};
  \draw (\rxcoordint,-0.45) -- (\rxcoordint,0.45);
  \node[anchor=north, inner sep=1.5pt] at (\rxcoordint,0) {\scriptsize $\rxcoordint$};
  \draw (-0.45,#2) -- (0.45,#2);
  \node[anchor=east, inner sep=1.5pt] at (0,#2) {\scriptsize $#2$};
\end{tikzpicture}}

\begin{document}

\begin{center}
{\large \textbf{Segunda Pr\'actica de Secundaria de la Prueba Nacional Estandarizada Sumativa de Matem\'atica 2026.}}
\end{center}

\noindent Nombre: \rule{8cm}{0.4pt}\hfill Sec: \rule{3cm}{0.4pt}

\vspace{6pt}
\noindent\textbf{INFORMACI\'ON GENERAL}

\vspace{4pt}
\noindent\textbf{Materiales necesarios para la prueba:}

\begin{itemize}[leftmargin=1.6em, itemsep=1pt, topsep=2pt, parsep=0pt]
  \item Un cuadernillo que contiene:
  \begin{itemize}[leftmargin=1.6em, itemsep=0pt, topsep=1pt, parsep=0pt, label=$\blacklozenge$]
    \item informaci\'on general
    \item 60 \'items de selecci\'on de respuesta
  \end{itemize}
  \item Hoja de respuestas para lectora \'optica
  \item Bol\'igrafo con tinta azul o negra
  \item Corrector l\'iquido blanco
\end{itemize}

\noindent\textbf{Material que se puede requerir para la prueba:}

\begin{itemize}[leftmargin=1.6em, itemsep=1pt, topsep=2pt, parsep=0pt]
  \item Calculadora
\end{itemize}

\vspace{4pt}
\noindent\textbf{Instrucciones:}

\begin{enumerate}[leftmargin=*, itemsep=3pt, topsep=2pt, parsep=0pt]
  \item La Prueba Nacional Estandarizada de secundaria de matem\'atica est\'a compuesta por 60 \'items. Verifique que el cuadernillo que tiene en sus manos est\'e bien compaginado y contenga los 60 \'items de selecci\'on de respuesta correspondientes al componente Matem\'aticas. En caso de encontrar alguna irregularidad, notif\'iquela inmediatamente al delegado de aula; de lo contrario, usted asume la responsabilidad sobre los problemas que se pudieran suscitar por esta causa.

  \item Cada \'item presenta tres posibilidades de respuesta: A), B), C) y D) Solamente una de ellas es la respuesta correcta.

  \item Lea cuidadosamente cada \'item y sus respectivas opciones. Puede utilizar el espacio al lado de cada \'item para realizar cualquier anotaci\'on que necesite, con el fin de hallar la respuesta.

  \item Ning\'un \'item debe aparecer sin respuesta o con m\'as de una marca en la hoja lectora \'optica.

  \item Una vez que haya revisado todas las opciones y tenga seguridad de su elecci\'on, rellene completamente el c\'irculo correspondiente, en la hoja lectora \'optica, tal como se indica en el siguiente ejemplo:

  \begin{center}
  \begin{tikzpicture}[baseline=-0.5ex]
    \draw[thick] (0,0) circle (0.38);
    \node at (0,0) {\textbf{A}};
    \fill (1.1,0) circle (0.38);
    \node[white] at (1.1,0) {\textbf{B}};
    \draw[thick] (2.2,0) circle (0.38);
    \node at (2.2,0) {\textbf{C}};
  \end{tikzpicture}
  \end{center}

  \item Si necesita rectificar la respuesta, utilice corrector l\'iquido blanco sobre el c\'irculo por corregir y rellene con bol\'igrafo de tinta negra o azul la nueva opci\'on seleccionada. Adem\'as, en el espacio de observaciones de la hoja lectora \'optica debe anotar y firmar la correcci\'on efectuada (Ejemplo: 12=A, firma). Se firma solo una vez al final de todas las correcciones.
\end{enumerate}

\vspace{8pt}
\noindent Creado por: Alberto Jim\'enez Ruiz asesor regional de Matem\'atica de la DRE-Cartago 2026.

\newpage

\begin{enumerate}[label=\textbf{\arabic*.}, itemsep=2.5cm]

\Needspace{0.92\textheight}
\item
En la plaza de un barrio existe una piscina comunal de forma circular.

La siguiente representaci\'on gr\'afica, cuyas unidades est\'an en metros, muestra las posiciones, $M$ del punto de referencia de la plaza, $D$ del centro de la piscina y $C$ de la circunferencia correspondiente al borde de esa piscina:

\begin{center}
\begin{tikzpicture}[scale=0.4]
  \path[use as bounding box] (-4,-7) rectangle (15,12);
  \draw[<->, thick] (-3,0) -- (13,0);
  \draw[<->, thick] (0,-6) -- (0,11);
  \node[anchor=west, inner sep=1pt] at (13.05,0) {\rotuloxeste};
  \node[anchor=south, inner sep=1pt] at (0,11.35) {\rotuloynorte};
  \draw[thick] (4,3) circle (6);
  \draw[dashed] (4,0) -- (4,3);
  \draw[dashed] (10,0) -- (10,3);
  \draw[dashed] (0,3) -- (10,3);
  \fill (0,0) circle (0.18) node[below left] {$M$};
  \fill (4,3) circle (0.18) node[anchor=south east, inner sep=1pt] {$D$};
  \fill (10,3) circle (0.18) node[anchor=south west, inner sep=1pt] {$C$};
  \node[anchor=north] at (4,0) {$4$};
  \node[anchor=north] at (10,0) {$10$};
  \node[left] at (0,3) {$3$};
\end{tikzpicture}
\end{center}

De acuerdo con la informaci\'on anterior, \textquestiondown cu\'al es la representaci\'on algebraica, cuyas unidades est\'an en metros, de la circunferencia $C$ que delimita esa piscina?

\begin{enumerate}[label=\Alph*)]
  \opcion{A}{(x-3)^2+(y-4)^2=36}
  \opcion{B}{(x+4)^2+(y+3)^2=12}
  \opcion{C}{(x-4)^2+(y-3)^2=12}
  \opcion{D}{(x-4)^2+(y-3)^2=36}
\end{enumerate}

\Needspace{0.90\textheight}
\item
En el jard\'in de una casa, se instala un aspersor de agua.

La siguiente representaci\'on gr\'afica, cuyas unidades est\'an en metros, muestra las posiciones, $V$ de la esquina del jard\'in, $A$ del aspersor y $C$ de la circunferencia correspondiente al alcance m\'aximo del chorro de agua lanzado por ese aspersor en metros:

\begin{center}
\begin{tikzpicture}[scale=0.4]
  \path[use as bounding box] (-9.5,-5) rectangle (8.2,11);
  \draw[<->, thick] (-8.5,0) -- (6,0);
  \draw[<->, thick] (0,-4) -- (0,10);
  \node[anchor=west, inner sep=1pt] at (6.05,0) {\rotuloxeste};
  \node[anchor=south, inner sep=1pt] at (0,10.35) {\rotuloynorte};
  \draw[thick] (-2,3) circle (4);
  \draw[dashed] (-2,0) -- (-2,3);
  \draw[dashed] (0,3) -- (-2,3);
  \draw[dashed] (2,0) -- (2,3);
  \draw[dashed] (0,3) -- (2,3);
  \fill (0,0) circle (0.18) node[below left] {$V$};
  \fill (-2,3) circle (0.18) node[anchor=south east, inner sep=1pt] {$A$};
  \fill (2,3) circle (0.18) node[anchor=south west, inner sep=1pt] {$C$};
  \node[anchor=north] at (-2,0) {$-2$};
  \node[anchor=north] at (2,0) {$2$};
  \node[left] at (0,3) {$3$};
\end{tikzpicture}
\end{center}

De acuerdo con la informaci\'on anterior, \textquestiondown cu\'al es la representaci\'on algebraica, cuyas unidades est\'an en metros, de la circunferencia $C$ correspondiente al alcance m\'aximo del chorro de agua?

\begin{enumerate}[label=\Alph*)]
  \opcion{A}{(x-2)^2+(y+3)^2=16}
  \opcion{B}{(x+2)^2+(y-3)^2=16}
  \opcion{C}{(x+2)^2+(y-3)^2=8}
  \opcion{D}{(x+2)^2+(y+3)^2=16}
\end{enumerate}

\Needspace{0.88\textheight}
\item
Una empresa instala un enrutador de Wi-Fi en el punto $E$ del plano de la oficina, donde las unidades est\'an en metros. El enrutador emite una se\~nal de $10$ metros a la redonda.

Cuatro computadoras est\'an ubicadas en los puntos $P(11,2)$, $Q(3,14)$, $R(-3,9)$ y $S(8,9)$.

La siguiente representaci\'on gr\'afica muestra la ubicaci\'on del enrutador y la circunferencia correspondiente al alcance m\'aximo de la se\~nal:

\begin{center}
\begin{tikzpicture}[scale=0.18]
  \path[use as bounding box] (-11,-12) rectangle (18.5,17);
  \draw[<->, thick] (-10,0) -- (16,0);
  \draw[<->, thick] (0,-11) -- (0,16);
  \node[anchor=west, inner sep=1pt] at (16.05,0) {\small\rotuloxeste};
  \node[anchor=south, inner sep=1pt] at (0,16.35) {\small\rotuloynorte};
  \draw[thick] (3,2) circle (10);
  \draw[dashed] (3,0) -- (3,2);
  \draw[dashed] (0,2) -- (3,2);
  \fill (3,2) circle (\puntografico);
  \node[anchor=south, inner sep=1pt, yshift=3pt] at (3,2) {$E$};
  \node[anchor=north] at (3,0) {$3$};
  \node[left] at (0,2) {$2$};
\end{tikzpicture}
\end{center}

De acuerdo con la informaci\'on anterior, \textquestiondown cu\'al de esas computadoras se debe reubicar para que tenga se\~nal Wi-Fi del enrutador?

\begin{enumerate}[label=\Alph*)]
  \item $Q$
  \item $S$
  \item $P$
  \item $R$
\end{enumerate}

\Needspace{0.88\textheight}
\item
Una empresa de telecomunicaciones modela el alcance de la se\~nal de una antena con la siguiente representaci\'on algebraica, donde las unidades est\'an en kil\'ometros:

\[(x+3)^2+(y-4)^2=9\]

\textquestiondown Cu\'al de las siguientes representaciones gr\'aficas corresponde a esa circunferencia?

\vspace{1\baselineskip}

\begin{center}
\setlength{\tabcolsep}{0pt}
\begin{tabular}{r@{\hspace{0.3cm}}c@{\hspace{0.6cm}}r@{\hspace{0.3cm}}c}
\raisebox{3.2cm}{\textbf{A)}} & \grafcirc{3}{4} & \raisebox{3.2cm}{\textbf{B)}} & \grafcirc{-3}{-4} \\[0.5cm]
\raisebox{3.2cm}{\textbf{C)}} & \grafcirc{-3}{4} & \raisebox{3.2cm}{\textbf{D)}} & \grafcirc{3}{-4}
\end{tabular}
\end{center}


\Needspace{0.72\textheight}
\item
Un sistema de alarma instalado en una bodega emite una se\~nal sonora circular. El alcance m\'aximo de esa se\~nal se modela mediante la circunferencia $C_1$ con la ecuaci\'on:
\[(x+4)^2+(y-2)^2=9\]
donde las unidades est\'an en metros.

Por razones de mantenimiento, el equipo se traslada $7$ metros hacia el este y $3$ metros hacia el sur.

De acuerdo con la informaci\'on anterior, \textquestiondown cu\'al es la ecuaci\'on de la circunferencia $C_2$ que representa el nuevo alcance de la se\~nal sonora?

\begin{enumerate}[label=\Alph*)]
  \opcion{A}{(x-3)^2+(y+1)^2=9}
  \opcion{B}{(x+3)^2+(y-1)^2=9}
  \opcion{C}{(x+11)^2+(y-5)^2=9}
  \opcion{D}{(x-3)^2+(y+1)^2=3}
\end{enumerate}

\Needspace{0.90\textheight}
\item
Un tel\'efono celular que emite una se\~nal inal\'ambrica circular se ubica en el punto $T$ de un dormitorio.

La siguiente representaci\'on gr\'afica, cuyas unidades est\'an en metros, muestra las ubicaciones $T$ del tel\'efono y $C$ de la circunferencia correspondiente al alcance m\'aximo de la se\~nal emitida:

\begin{center}
\begin{tikzpicture}[scale=0.45]
  \path[use as bounding box] (-3,-2) rectangle (10,12);
  \draw[<->, thick] (-2.5,0) -- (8,0);
  \draw[<->, thick] (0,-1.5) -- (0,11);
  \node[anchor=west, inner sep=1pt] at (8.05,0) {\rotuloxeste};
  \node[anchor=south, inner sep=1pt] at (0,11.35) {\rotuloynorte};
  \draw[thick] (2,5) circle (3);
  \draw[dashed] (2,0) -- (2,5);
  \draw[dashed] (5,0) -- (5,5);
  \draw[dashed] (0,5) -- (5,5);
  \fill (2,5) circle (0.15) node[anchor=south east, inner sep=1pt] {$T$};
  \fill (5,5) circle (0.15) node[anchor=south west, inner sep=1pt] {$C$};
  \node[anchor=north] at (2,0) {$2$};
  \node[anchor=north] at (5,0) {$5$};
  \node[left] at (0,5) {$5$};
\end{tikzpicture}
\end{center}

De acuerdo con la informaci\'on anterior, si ese tel\'efono se traslada $4$ metros hacia la izquierda y $5$ metros hacia el sur, \textquestiondown cu\'al es la representaci\'on algebraica, cuyas unidades est\'an en metros, de la circunferencia correspondiente al alcance m\'aximo de la se\~nal en su nueva ubicaci\'on?

\begin{enumerate}[label=\Alph*)]
  \opcion{A}{(x-2)^2+y^2=9}
  \opcion{B}{(x+2)^2+y^2=9}
  \opcion{C}{(x+2)^2+y^2=3}
  \opcion{D}{(x+2)^2+(y-5)^2=9}
\end{enumerate}

\Needspace{0.92\textheight}
\item
Un aspersor de riego instalado en un terreno cubre un \'area circular cuyo alcance m\'aximo se modela mediante la circunferencia $C_1$ con ecuaci\'on:
\[(x-5)^2+(y+3)^2=16\]
donde las unidades est\'an en metros.

Debido a una remodelaci\'on del terreno, el aspersor se reubica $3$ metros hacia el oeste y $6$ metros hacia el norte.

De acuerdo con la informaci\'on anterior, \textquestiondown cu\'al de las siguientes representaciones gr\'aficas, cuyas unidades est\'an en metros, corresponde al alcance del aspersor en su nueva ubicaci\'on?

\vspace{1\baselineskip}

\begin{center}
\setlength{\tabcolsep}{0pt}
\begin{tabular}{r@{\hspace{0.4cm}}c@{\hspace{1.2cm}}r@{\hspace{0.4cm}}c}
\raisebox{3.5cm}{\textbf{A)}} & \grafcirctrasl{2}{3}{7} & \raisebox{3.5cm}{\textbf{B)}} & \grafcirctrasl{8}{3}{4} \\[1.8cm]
\raisebox{3.5cm}{\textbf{C)}} & \grafcirctrasl{2}{-9}{4} & \raisebox{3.5cm}{\textbf{D)}} & \grafcirctrasl{2}{3}{4}
\end{tabular}
\end{center}

\Needspace{0.90\textheight}
\item
En una feria cient\'ifica se coloca un sensor circular sobre una mesa para detectar la posici\'on final de peque\~nas fichas met\'alicas. El borde de la zona de detecci\'on se representa mediante una circunferencia $C$ de centro $O$.

La figura muestra, en cent\'imetros, la circunferencia $C$ y su centro $O$:

\begin{center}
\begin{tikzpicture}[scale=0.45]
  \path[use as bounding box] (-2,-2) rectangle (13,11);
  \draw[<->, thick] (-1,0) -- (12,0);
  \draw[<->, thick] (0,-1) -- (0,10);
  \node[anchor=west, inner sep=1pt] at (12.05,0) {$x$};
  \node[anchor=south, inner sep=1pt] at (0,10.35) {$y$};
  \draw[thick] (7,5) circle (3);
  \draw[dashed] (7,0) -- (7,5);
  \draw[dashed] (4,0) -- (4,5);
  \draw[dashed] (0,5) -- (7,5);
  \fill (7,5) circle (0.15);
  \fill (4,5) circle (0.15);
  \node[anchor=south, inner sep=1pt, yshift=2pt] at (7,5) {$O$};
  \node[anchor=north] at (7,0) {$7$};
  \node[anchor=north] at (4,0) {$4$};
  \node[left] at (0,5) {$5$};
\end{tikzpicture}
\end{center}

Durante una prueba, cuatro fichas identificadas como $R$, $S$, $W$ y $T$ quedaron ubicadas, respectivamente, en los puntos $(9,7)$, $(4,5)$, $(10,9)$ y $(7,2)$.

De acuerdo con la informaci\'on anterior, \textquestiondown cu\'al de esas fichas qued\'o ubicada en el exterior de la zona de detecci\'on?

\begin{enumerate}[label=\Alph*)]
  \item $R$
  \item $S$
  \item $T$
  \item $W$
\end{enumerate}

\Needspace{0.96\textheight}
\item
Considere la siguiente informaci\'on:

La siguiente figura representa dos zonas circulares de jardiner\'ia, $C_1$ de centro $O$ y radio $5$, y $C_2$ de centro $P$ y radio $3$, adem\'as de la recta $m$, en un sistema de coordenadas:

\begin{center}
\begin{tikzpicture}[scale=0.48]
  \path[use as bounding box] (-7.2,-1.7) rectangle (14.2,12.2);
  \draw[<->, thick] (-6,0) -- (13,0);
  \draw[<->, thick] (0,-1) -- (0,11);
  \node[anchor=west, inner sep=1pt] at (13.05,0) {$x$};
  \node[anchor=south, inner sep=1pt] at (0,11.25) {$y$};
  \draw[thick] (0,5) circle (5);
  \draw[thick] (8,3) circle (3);
  \draw[very thick, <->] (-6,6) -- (13,6) node[anchor=south west, inner sep=1pt] {$m$};
  \draw[dashed] (8,0) -- (8,3);
  \draw[dashed] (0,3) -- (8,3);
  \fill (0,5) circle (0.15);
  \fill (8,3) circle (0.15);
  \fill (0,0) circle (0.15);
  \fill (8,0) circle (0.15);
  \fill (8,6) circle (0.15);
  \fill (-4.9,6) circle (0.15);
  \node[anchor=south west, inner sep=1pt, xshift=2pt] at (0,5) {$O$};
  \node[anchor=south west, inner sep=1pt] at (8,3) {$P$};
  \node[anchor=north east, inner sep=1pt] at (0,0) {$A$};
  \node[anchor=south east, inner sep=1pt] at (8,0) {$B$};
  \node[anchor=south west, inner sep=1pt] at (8,6) {$D$};
  \node[anchor=south east, inner sep=1pt] at (-4.9,6) {$E$};
  \node[anchor=north, yshift=-6pt] at (8,0) {$8$};
  \node[left] at (0,3) {$3$};
  \node[left] at (0,5) {$5$};
  \node[anchor=south west, inner sep=1pt] at (4.3,9.2) {$C_1$};
  \node[anchor=west, inner sep=1pt] at (11.2,3.7) {$C_2$};
\end{tikzpicture}
\end{center}

De acuerdo con la informaci\'on de la figura anterior, considere las siguientes proposiciones:

I. El eje $x$ es tangente a $C_1$ y a $C_2$.

II. La recta $m$ es secante a $C_1$ y tangente a $C_2$.

De ellas, \textquestiondown cu\'al o cu\'ales son verdaderas?

\begin{enumerate}[label=\Alph*)]
  \item Ambas.
  \item Ninguna.
  \item Solo la I.
  \item Solo la II.
\end{enumerate}

\Needspace{0.92\textheight}
\item
Considere la siguiente representaci\'on gr\'afica, en la cual la circunferencia $C$ de centro $M$ y la circunferencia $P$ de centro $H$ son tangentes en el punto $B$, y el eje $x$ es tangente a $C$ en el punto $A$:

\begin{center}
\begin{tikzpicture}[scale=0.95, >=stealth]
  \path[use as bounding box] (-1.7,-5.1) rectangle (9.3,4.4);
  \def\puntotangencia{0.12}
  \draw[<->, thick] (-1.4,0) -- (8.9,0);
  \draw[<->, thick] (0,-3.8) -- (0,3.9);
  \node[anchor=west, inner sep=1pt] at (8.95,0) {$x$};
  \node[anchor=south, inner sep=1pt] at (0,4.05) {$y$};
  \draw[thick] (0,-1.4) circle (1.4);
  \draw[thick] (4.8,0) circle (3.6);
  \draw[dashed, thick] (0,-1.4) -- (4.8,0);
  \fill (0,0) circle (\puntotangencia) node[anchor=south east, inner sep=1pt, xshift=-2pt, yshift=2pt] {$A$};
  \fill (0,-1.4) circle (\puntotangencia) node[anchor=north east, inner sep=1pt] {$M$};
  \fill (4.8,0) circle (\puntotangencia) node[anchor=south, inner sep=1pt, yshift=3pt] {$H$};
  \fill (1.34,-1.01) circle (\puntotangencia) node[anchor=south west, inner sep=1pt, xshift=2pt, yshift=1pt] {$B$};
  \node[anchor=south west, inner sep=1pt] at (-0.9,-2.4) {$C$};
  \node[anchor=south, inner sep=1pt] at (4.8,3.6) {$P$};
\end{tikzpicture}
\end{center}

De acuerdo con la informaci\'on anterior, si $AH=48$ y $MH=50$, entonces, \textquestiondown cu\'al es la medida del radio de $P$?

\begin{enumerate}[label=\Alph*)]
  \opcion{A}{14}
  \opcion{B}{36}
  \opcion{C}{50}
  \opcion{D}{64}
\end{enumerate}

\Needspace{0.82\textheight}
\item
La gr\'afica adjunta representa un camino $m$ construido en un terreno cuyas unidades est\'an en kil\'ometros:

\begin{center}
\begin{tikzpicture}[scale=0.62]
  \path[use as bounding box] (-2.6,-2.2) rectangle (6.5,5.3);
  \draw[<->, thick] (-2.2,0) -- (5.8,0);
  \draw[<->, thick] (0,-1.8) -- (0,4.8);
  \node[anchor=west, inner sep=1pt] at (5.85,0) {$x$};
  \node[anchor=south, inner sep=1pt] at (0,4.95) {$y$};
  \draw[very thick, <->] (-1.4,4.05) -- (5.35,-1.01);
  \node[anchor=south east, inner sep=1pt] at (-1.15,4.05) {$m$};
  \fill (0,3) circle (\puntografico);
  \fill (4,0) circle (\puntografico);
  \node[anchor=east] at (0,3) {$3$};
  \node[anchor=north] at (4,0) {$4$};
\end{tikzpicture}
\end{center}

Si se desea crear un camino perpendicular al camino $m$ representado por $3y=kx-6$, entonces, \textquestiondown cu\'al es el valor de $k$?

\begin{enumerate}[label=\Alph*)]
  \opcion{A}{4}
  \opcion{B}{\frac{3}{4}}
  \opcion{C}{-4}
  \opcion{D}{-\frac{3}{4}}
\end{enumerate}

\Needspace{0.88\textheight}
\item
A continuaci\'on, se muestra la representaci\'on gr\'afica, cuyas unidades est\'an en metros, del piso de una casa:

\begin{center}
\begin{tikzpicture}[scale=0.42]
  \path[use as bounding box] (-1.2,-1.1) rectangle (13.8,9.6);
  \draw[->, thick] (-0.8,0) -- (13.2,0);
  \draw[->, thick] (0,-0.8) -- (0,9.0);
  \node[anchor=west, inner sep=1pt] at (13.25,0) {$x$};
  \node[anchor=south, inner sep=1pt] at (0,9.1) {$y$};
  \draw[very thick] (0,8) -- (12,8) -- (12,0) -- (4,0) -- (0,4) -- cycle;
  \draw[thick] (0.45,8) -- (0.45,7.55) -- (0,7.55);
  \draw[thick] (11.35,0) -- (11.35,0.65) -- (12,0.65);
  \fill (0,8) circle (\puntografico);
  \fill (0,4) circle (\puntografico);
  \fill (4,0) circle (\puntografico);
  \fill (12,0) circle (\puntografico);
  \node[anchor=east] at (0,8) {$8$};
  \node[anchor=east] at (0,4) {$4$};
  \node[anchor=north] at (4,0) {$4$};
  \node[anchor=north] at (12,0) {$12$};
\end{tikzpicture}
\end{center}

De acuerdo con la informaci\'on anterior, si se le quiere colocar cer\'amica a la totalidad del piso de esa casa, entonces, \textquestiondown cu\'antos metros cuadrados de cer\'amica, como m\'inimo, se deben comprar?

\begin{enumerate}[label=\Alph*)]
  \opcion{A}{80}
  \opcion{B}{88}
  \opcion{C}{96}
  \opcion{D}{104}
\end{enumerate}

\Needspace{0.78\textheight}
\item
En un centro recreativo se construir\'an dos plazoletas. La primera tendr\'a forma de dodec\'agono regular y la segunda forma de oct\'agono regular.

La medida del lado del dodec\'agono debe ser el doble de la medida del lado del oct\'agono.

Si el per\'imetro de la plazoleta con forma de dodec\'agono debe ser $96\ \text{m}$, entonces, \textquestiondown cu\'al ser\'a el per\'imetro de la plazoleta con forma de oct\'agono?

\begin{enumerate}[label=\Alph*)]
  \opcion{A}{32\ \text{m}}
  \opcion{B}{48\ \text{m}}
  \opcion{C}{64\ \text{m}}
  \opcion{D}{96\ \text{m}}
\end{enumerate}

\Needspace{0.92\textheight}
\item
La siguiente representaci\'on gr\'afica muestra la superficie de una pieza decorativa que se colocar\'a sobre una pared:

\begin{center}
\fbox{\begin{minipage}{0.42\textwidth}
\centering
La medida del lado de cada cuadrado de la cuadr\'icula es $10\ \text{cm}$.
\end{minipage}}

\vspace{0.2cm}

\begin{tikzpicture}[scale=0.72]
  \path[use as bounding box] (-0.8,-0.8) rectangle (9.2,8.2);
  \draw[step=1, dotted, black!68] (0,0) grid (8,7);
  \draw[<->, thick] (-0.4,0) -- (8.6,0);
  \draw[<->, thick] (0,-0.4) -- (0,7.6);
  \node[anchor=west, inner sep=1pt] at (8.65,0) {$x$};
  \node[anchor=south, inner sep=1pt] at (0,7.75) {$y$};
  \draw[very thick] (2,1) -- (2,6) -- (6,4) -- (6,2) -- cycle;
  \fill (2,1) circle (\puntografico);
  \fill (2,6) circle (\puntografico);
  \fill (6,4) circle (\puntografico);
  \fill (6,2) circle (\puntografico);
\end{tikzpicture}
\end{center}

De acuerdo con la informaci\'on anterior, si se requiere colocar una cinta adhesiva alrededor de todo el borde de esa pieza decorativa, entonces, \textquestiondown cu\'al es la cantidad m\'inima aproximada de cent\'imetros de cinta que se requiere?

\begin{enumerate}[label=\Alph*)]
  \opcion{A}{140}
  \opcion{B}{156}
  \opcion{C}{180}
  \opcion{D}{200}
\end{enumerate}

\Needspace{0.88\textheight}
\item
Una carpintera dise\~n\'o una puerta decorativa como la que se muestra en la siguiente representaci\'on gr\'afica:

\begin{center}
\fbox{\begin{minipage}{0.42\textwidth}
\centering
La medida del lado de cada cuadrado de la cuadr\'icula es $1\ \text{m}$.
\end{minipage}}

\vspace{0.55cm}

\begin{tikzpicture}[scale=0.62]
  \path[use as bounding box] (-0.8,-0.8) rectangle (6.8,6.4);
  \draw[step=1, dotted, gray!85] (0,0) grid (6,6);
  \draw[<->, thick] (-0.4,0) -- (6.4,0);
  \draw[<->, thick] (0,-0.4) -- (0,6.4);
  \node[anchor=west, inner sep=1pt] at (6.45,0) {$x$};
  \node[anchor=south, inner sep=1pt] at (0,6.55) {$y$};
  \draw[very thick] (1,1) -- (5,1);
  \draw[very thick] (1,1) -- (1,3);
  \draw[very thick] (5,1) -- (5,3);
  \draw[very thick] (1,3) arc[start angle=180,end angle=0,radius=2];
\end{tikzpicture}
\end{center}

De acuerdo con la informaci\'on anterior, el per\'imetro, en metros, de esa puerta decorativa es

\begin{enumerate}[label=\Alph*)]
  \item mayor que $12$, pero menor que $13$.
  \item mayor que $14$, pero menor que $15$.
  \item mayor que $15$, pero menor que $16$.
  \item mayor que $10$, pero menor que $12$.
\end{enumerate}

\end{enumerate}

\end{document}
